%قبل از شروع مسئله ترجیحاً چیزی ننویسید. اما اگر خواستید بنویسید، در قالب لیستی (فقرات و شمارش) بنویسید.
%برای اضافه کردن نمره، بعدی از شروع‌مسئله، [x] را اضافه کنید و به جای x  نمره مورد نظر را قرار دهید. برای زیرمسئله‌ها هم به همین صورت است.
%\شروع‌مسئله[10]
\شروع‌مسئله [10+10]
% صورت اصلی مسئله را در این‌جا بنویسید.


\شروع{زیرمسائل}
\زیرمسئله
 الگوریتم مرتب سازی زیر را  در نظر بگیرید (این الگوریتم مرتب سازی به \متن‌لاتین{selection sort} معروف است)

\begin{latin}
\begin{lstlisting}
def min_index(a, i, j):
	if i == j:
		return i
	k = min_index(a, i + 1, j)

	if a[i] < a[k]:
		return i
	else:
		return k


def selection_sort(a, n, index):
	if index == n:
		return -1
	k = min_index(a, index, n - 1)

	if k != index:
		a[k], a[index] = a[index], a[k]
	selection_sort(a, n, index + 1)


def main():
	arr = [32, 11, -5, 1, 7, 61, 43, 3]
	selection_sort(arr, len(arr), 0)
	
main()
\end{lstlisting}
\end{latin}


\شروع{شمارش}
\فقره درستی این الگوریتم را ثابت کنید.
\فقره زمان اجرای آن را تحلیل کنید.
%\زیرمسئله % مسئلۀ ۱پ
%\زیرمسئله % مسئلۀ ۱ت
\پایان{شمارش}



\زیرمسئله
یک مدل از مرتب‌سازی \متن‌لاتین{merge sort} برای یک آرایه $A$ به صورت زیر می‌باشد:(منظور از 
$A[i..j]$
بخشی از آرایه که بین 
$A[i]$
و
$A[j]$
هست می‌باشد
)

\شروع{فقرات}
\فقره اگر 
$n=1$
در آن صورت الگوریتم به پایان می‌رسد.
\فقره در غیر اینصورت قرار دهید
$k=\ceil{\frac{n}{3}}$
\فقره به صورت بازگشتی، 
$A[1..k]$
و
$A[k+1..n]$
را مرتب کنید.
\فقره دو زیرآرایه 
$A[1..k]$
و
$A[k+1..n]$
را ادغام کنید.

\پایان{فقرات}

ثابت کنید پیچیدگی زمانی این الگوریتم برابر است با 
$\mathcal{O}(n\log{}n)$
.
\پایان{زیرمسائل}
%اگر مسئله چند قسمت داشت به این صورت عمل کنید:


\solution{
\شروع{فقرات}
\فقره
	در واقع این الگوریتم در مرحله $k$اُم عدد $k$اُم را (از لحاظ کوچکی) سر جای خودش قرار می دهد پس بعد از $n$ مرحله هر عددی سر جای خودش قرار می‌گیرد. (در مرحله اول عنصر کمینه را در جایگاه اول قرار می دهد سپس در گام بعدی عنصر کمینه بعدی را در جایگاه دوم قرار می دهد و ... تا به آخر آرایه برسد)
	
به ازای هر حالت اولیه‌ای در زمان
$O(n^2)$
آرایه را مرتب
می‌کند. چون هر دفعه باید از بین $k$ عدد، کمینه را پیدا کند
($1 \leq k \leq n$)
پس تعداد عملیات انجام شده برابر است با
$1 + 2 + ... + n = O(n^2)$.

\فقره
هر دفعه آرایه به دو قسمت 
$\ceil{\frac{n}{3}}$تایی و 
$\floor{\frac{2n}{3}}$تایی 
تقسیم می‌شود. پس رابطه بازگشتی زمان اجرای این الگوریتم برابر است با
$T(n) = T(\ceil{\frac{n}{3}}) + T(\floor{\frac{2n}{3}}) + O(n)$.
(توجه کنید ادغام کردن مانند ادغام 
\متن‌لاتین{merge sort}
معمولی در 
$O(n)$
انجام می شود)

با ترسیم درخت بازگشت برای این رابطه بازگشتی متوجه می شویم که جواب آن
$O(n\log n)$
است چون تعداد طبقات 
$\log_{1/3} n$
است و در هر طبقه عملیات در مجموع از مرتبه 
$O(n)$
است.

\پایان{فقرات}
}




%%%%%%%%%%%%%%%%%%%%%%%%%%%%%%%%%%%%%%%%%%%%%%%%%%%%%
% برخی دستورات مفید برای تمرین‌ها. %
%%%%%%%%%%%%%%%%%%%%%%%%%%%%%%%%%%%%%%%%%%%%%%%%%%%%%

% فرمول‌های نمونه:
%$f(x) = \Theta(x)$
%$T(x, y) \leq \log(x) + 2^y + \binom{2n}{n}$
% {\tt code\_function}


% لیست موارد بدون شماره‌گذاری:
%\شروع{فقرات}
%    \فقره مورد اول 
%        \شروع{فقرات}
%            \فقره مورد اول داخلی 
%                \شروع{فقرات}
%                    \فقره اولین مورد خیلی داخلی 
%                    \فقره دومین مورد خیلی داخلی 
%                \پایان{فقرات}
%            \فقره مورد دوم داخلی 
%        \پایان{فقرات}
%    \فقره مورد دوم 
%\پایان{فقرات}

% لیست موارد شماره‌گذاری شده:
%\شروع{شمارش}
%    \فقره یک مورد 
%    \فقره یک مورد دیگر 
%    \فقره باز هم یک مورد دیگر
%\پایان{شمارش}

% برای نوشتن یک فرمول چند خطی که یک شماره بخوره. برای حالت‌های دیگر به سایت زیر مراجعه کنید:
% https://www.overleaf.com/learn/latex/Aligning_equations_with_amsmath
%\begin{align} 
%    \begin{split}
%        (x+y)^3 &= (x+y)^2(x+y) \\
%                &= (x^2+2xy+y^2)(x+y) \\
%                &= (x^3+2x^2y+xy^2) + (x^2y+2xy^2+y^3) \\
%                &= x^3+3x^2y+3xy^2+y^3
%    \end{split}                                 
%\end{align}

% نوشتن ماتریس و آرایه دوبعدی
%\begin{align}
%    A = 
%    \begin{bmatrix}
%        A_{11} & A_{21} \\
%        A_{21} & A_{22}
%    \end{bmatrix}
%\end{align}

% اضافه کردن تصویر با شرح. برای اضافه کردن تصویر بدون شرح، فقط دستور درج‌تصویر را استفاده کنید.
%\شروع{شکل}[h!]
%\تنظیم‌ازوسط
%  \درج‌تصویر[پهنا=0.5\پهنای‌متن]{img.jpg}
%\شرح{شرح تصویر}
%\برچسب{تصویر:۱}
%\پایان{شکل}
%\پایان{زیرمسائل}

%کد زیر را ببینید.
%%\lstinputlisting{source_filename.py}
%\begin{latin}
%\begin{lstlisting}
%# Import the modules
%import sys
%import random
%
%ans = True
%
%while ans:
%question = raw_input("Ask the magic 8 ball a question: (press enter to quit) ")
%
%answers = random.randint(1,8)
%
%if question == "":
%sys.exit()
%
%elif answers == 1:
%print "It is certain"
%
%elif answers == 2:
%print "Outlook good"
%
%elif answers == 3:
%print "You may rely on it"
%
%elif answers == 4:
%print "Ask again later"
%
%elif answers == 5:
%print "Concentrate and ask again"
%
%elif answers == 6:
%print "Reply hazy, try again"
%
%elif answers == 7:
%print "My reply is no"
%
%elif answers == 8:
%print "My sources say no"
%\end{lstlisting}
%\end{latin}
